% Each \section is a new appendix (A, B, ...).
% Figures are numbered within their appendix (A1, A2, ..., B1, ...),
% so put them in the appendix section they belong to.
% Put a \clearpage before each \section and \subsection and at the end of the file,
% so each appendix's figures are flushed before the next one starts.
% The per-gas figures (the `@...-appendix-figures@` lines below) are generated by the build
% (`src/local/scenario_ghg_forcing_for_cmip7/appendix_figures.py`), so edit their captions there.
% A figure which is already included anywhere in the text is left out of the generated figures,
% so it isn't in the manuscript twice.
% \appendix and \noappendix are added by the build, don't add them here.
% See `manuscripts/historical-ghg-forcing-for-cmip7/appendices.tex` for why these are appendices
% (and not a supplement) and why we don't use \appendixfigures.

\clearpage
\section{Additional methods figures}
\label{app:methods-figures}

\subsection{Regressions}
\label{app:methods-figures-regressions}

% Checked 2026-10-26: the CH_4 regression shown is the one used for the published data
% (the slope implied by the published L data, -0.74943 yr / Mt CH_4, matches the re-derived regression exactly).
% It uses all historical years in which the first PC isn't constant (1750-2023),
% so it mixes the period where the historical first PC comes from the optimisation against ice cores (before 1948)
% with the period where it comes from the regression against emissions.
% Its slope is also around 5.6 times that of the historical dataset's regression against PRIMAP-hist (-0.135 yr / Mt CH_4).
% ZNTODO: wherever we refer to this regression (here and in Sect. \ref{sssec:methods-ch4-project-seasonality-lat-gradient}),
% refer to the discussion of inconsistencies between inputs (Sect. \ref{sec:discussion}).

\begin{figure}[hbt!]
\includegraphics[width=12cm]{@regressions-methods-figure@}
\caption{
    Regressions we re-derive for the projections
    (Sects.~\ref{sssec:methods-co2-project-seasonality-lat-gradient},
    \ref{sssec:methods-ch4-project-seasonality-lat-gradient}
    and~\ref{sssec:methods-n2o-like-project-seasonality-lat-gradient}).
    a)--c) First latitudinal gradient PC from the historical dataset
    against the historical emissions which underpin the scenarios:
    a) fossil CO_2, b) fossil CH_4, c) total N_2O.
    d) CO_2 seasonality change PC from the historical dataset against MAGICC's median net primary productivity (NPP).
    Points are coloured by year, the black lines are the regressions.
    For the other gases, we use the regressions from the historical dataset [TODO cite historical paper] directly,
    so they are not shown here.
}
\label{fig:methods-regressions}
\end{figure}

\clearpage
\subsection{Extension of the Western et al. (2024) projections}
\label{app:methods-figures-western-et-al-2024}

\begin{figure}[hbt!]
\includegraphics[width=12cm]{@western-et-al-2024-extensions-figure@}
\caption{
    Extension of the \citet{western_hcfcs_2024} projections beyond 2100
    (Sect.~\ref{sssec:methods-cfc12-like-project-global-mean}).
    a) HCFC-22, b) HCFC-141b, c) HCFC-142b.
    In each panel: the historical dataset, the \citet{western_hcfcs_2024} projection (to 2100),
    our extension of it (after 2100)
    and the \citet{wmo_ozone-ch7_2022} projection, for reference.
}
\label{fig:methods-western-et-al-2024}
\end{figure}

\clearpage
\subsection{Gases processed like N_2O}
\label{app:methods-figures-n2o-like}

@n2o-like-methods-appendix-figures@

\clearpage
\subsection{Gases processed like CFC-12}
\label{app:methods-figures-cfc12-like}

@cfc12-like-methods-appendix-figures@

\clearpage
\subsection{Gases processed like HFC-125}
\label{app:methods-figures-hfc125-like}

@hfc125-like-methods-appendix-figures@

\clearpage
\section{MAGICC compared with our concentrations for the gases which use a single concentration projection}
\label{app:magicc-mismatch}

For the gases which use a single concentration projection,
MAGICC is run with the official emissions, including the official extensions beyond 2100,
which differ from the emissions implied by our concentrations after 2100
(Sect.~\ref{ssec:methods-input-emissions}).
The figures in this appendix show the size of this difference for each gas.

@magicc-mismatch-appendix-figures@

\clearpage
